\originalchapter{角动量与中心势}
\originalsection{角动量代数}
中心力沿半径, 经典力矩$\boldsymbol r\times\boldsymbol F=0$; 量子中心势也不选空间方向, 旋转不改变Hamiltonian. 因而应寻找旋转生成的守恒量. $L^2$测总轨道角动量大小的平方, $L_z$测相对选定轴的投影; 不同分量不对易, 不能同时赋确定$x,y,z$读数. $\boldsymbol L$为轨道,$\boldsymbol S$为内禀自旋,$\boldsymbol J$可为任一种或它们之和. $j,m$无量纲, 实际读数带$\hbar$单位. 轨道角动量定义为$\boldsymbol L=\boldsymbol r\times\boldsymbol p$, $\boldsymbol p=-\ii\hbar\nabla$. 在光滑函数的共同定义域, 利用$[x_i,p_j]=\ii\hbar\delta_{ij}$逐项展开, 可得$[L_x,L_y]=\ii\hbar L_z$及循环式. $L^2=L_x^2+L_y^2+L_z^2$与三个分量都对易. 因而可选$L^2,L_z$共同本征态. 以下对一般角动量$\boldsymbol J$使用同一代数.

设$J^2\ket{j,m}=\hbar^2j(j+1)\ket{j,m}$, $J_z\ket{j,m}=\hbar m\ket{j,m}$. 初始可把$J^2$本征值写为非负参数$\hbar^2\lambda$; 下面的终止条件说明$\lambda=j(j+1)$的来源. 定义$J_\pm=J_x\pm\ii J_y$, 则$[J_z,J_\pm]=\pm\hbar J_\pm$且$[J^2,J_\pm]=0$. 升降只改变$m$, 不改变$J^2$.

由 $J_-J_+=J^2-J_z^2-\hbar J_z$ 和 $J_+J_-=J^2-J_z^2+\hbar J_z$, 范数平方为$\hbar^2[\lambda-m(m\pm1)]$. 它必须非负, 所以$m$不能无限升降. 若最高和最低值为$m_+,m_-$, 终止给$\lambda=m_+(m_++1)=m_-(m_--1)$. 相减因式分解得$(m_+-m_-+1)(m_++m_-)=0$, 前一因子正, 所以$m_-=-m_+$; “表示”在此指由升降相连的一组态, 不是粒子绕过多少圈. 记$m_+=j\ge0$. 两端之间的整数步数为$2j$, 因而$j$为非负整数或半整数, $m=-j,-j+1,\ldots,j$. 归一化系数为
\[
J_\pm\ket{j,m}=\hbar\sqrt{(j\mp m)(j\pm m+1)}\ket{j,m\pm1}.
\]
这是角动量量子数和系数的代数推导, 不能仅从“整数轨道圈数”推出自旋半整数.

\originalsection{球谐函数与径向方程}
球坐标定义为$x=r\sin\theta\cos\phi,y=r\sin\theta\sin\phi,z=r\cos\theta$. 中心势角方向相同, 可先找乘积$R(r)Y(\theta,\phi)$, 用角函数分离旋转性质. $L^2=-\hbar^2[\partial_\theta^2+\cot\theta\partial_\theta+\csc^2\theta\partial_\phi^2]$, 与$L_z$一起给角向本征方程. 在南北极正则、方位单值的解限制$\ell,m$, 本册不要求预先会全部关联Legendre函数. 以下$\mu$表示相对运动的约化质量, 固定中心单粒子时取粒子质量; 两体约化在下一章展开. 对于轨道角动量, 球坐标中$L_z=-\ii\hbar\partial_\phi$, 所以$L_z$本征函数的方位因子为$\e^{\ii m\phi}$. 位置波函数对$\phi\mapsto\phi+2\pi$单值, 强制$m$为整数. 因而轨道角动量的$j=\ell$也为整数. 角向本征函数称为球谐函数$Y_{\ell m}(\theta,\phi)$, 归一化为$\int|Y_{\ell m}|^2\dd\Omega=1$, $\dd\Omega=\sin\theta\dd\theta\dd\phi$. 最简单$Y_{00}=1/\sqrt{4\pi}$给每立体角等概率, 因$\int\dd\Omega=4\pi$而归一. $Y_{10}=\sqrt{3/(4\pi)}\cos\theta$, 代$L^2$得$2\hbar^2Y_{10}$, 代$L_z$得0; $\int\cos^2\theta\dd\Omega=4\pi/3$检验归一. 这分别是$s$与一个$p$角分布, $p$不是沿确定轴的经典轨道. 这与自旋量子数的半整数可能性并不冲突, 因为自旋不以这种单值标量角函数表示.

若$V=V(r)$, 三维方程分离为$\psi=R(r)Y_{\ell m}$. 使用
$\nabla^2=r^{-2}\partial_r(r^2\partial_r)-L^2/(\hbar^2r^2)$, 得
\[
-\frac{\hbar^2}{2\mu}\left(R''+\frac2rR'\right)
+\left[V(r)+\frac{\hbar^2\ell(\ell+1)}{2\mu r^2}\right]R=ER.
\]
取$u=rR$, 实际$R\prime=u\prime/r-u/r^2$, $R\prime\prime+2R\prime/r=u\prime\prime/r$, 因而一阶导数项消失:
\begin{equation}\label{eq:radial}
-\frac{\hbar^2}{2\mu}u''+
\left[V(r)+\frac{\hbar^2\ell(\ell+1)}{2\mu r^2}\right]u=Eu,
\qquad r>0.
\end{equation}
体积元$\dd^3r=r^2\dd r\dd\Omega$, 所以$\int|\psi|^2\dd^3r=\int|R|^2r^2\dd r\int|Y|^2\dd\Omega=\int|u|^2\dd r$. $|u|^2\dd r$是薄球壳概率, 不是某个空间点的密度. 归一化是$\int_0^\infty|u|^2\dd r=1$, 不是$\int|R|^2\dd r=1$. 对通常正则中心势, 物理边界为$u(0)=0$. 额外的离心项把三维问题化为半轴上的有效一维问题. 它不能与全实线的任意一维势阱混同.

\originalsection{角动量合成}
先把联合测量标签写清: $\ket{+-}$表示系统1读$+\hbar/2$、系统2读$-\hbar/2$的乘积制备. $\otimes$把两个态组合, 满足$(a u+bv)\otimes w=a(u\otimes w)+b(v\otimes w)$, 一侧算符只作用该侧. 四个自旋组合形成四维基, 第12章系统讨论一般联合态. 对两个独立系统, $\boldsymbol J=\boldsymbol J^{(1)}+\boldsymbol J^{(2)}$, 不同系统的分量对易, 所以总角动量仍满足同一代数. 乘积基$\ket{j_1,m_1}\otimes\ket{j_2,m_2}$使$J_z$本征值为$\hbar(m_1+m_2)$, 但一般不是$J^2$本征态.

两个自旋$1/2$的最高总磁量子数态是$\ket{++}$, $j=1,m=1$. $J_-=S_{1-}+S_{2-}$, 每个分量分别给$\hbar\ket{-+}$和$\hbar\ket{+-}$, 合态范数为$\sqrt2\hbar$. 用$J_-$降阶, 得三重态$\ket{1,0}=(\ket{+-}+\ket{-+})/\sqrt2$, $\ket{1,-1}=\ket{--}$. 剩下的正交态是$\ket{0,0}=(\ket{+-}-\ket{-+})/\sqrt2$, 直接验证被$J_z$和$J_\pm$湮灭. 四维空间分为三维$j=1$和一维$j=0$子空间.

又因为$J^2=(J^{(1)})^2+(J^{(2)})^2+2\boldsymbol J^{(1)}\cdot\boldsymbol J^{(2)}$, 若相互作用为$\lambda\boldsymbol J^{(1)}\cdot\boldsymbol J^{(2)}$, 两类态的能量移分别为$\lambda\hbar^2/4$与$-3\lambda\hbar^2/4$. 这给自旋耦合的分裂, 不需要对原来的$4\times4$矩阵盲目求特征多项式.

\originalsection{轨道与自旋二分之一的合成}
为下一章分裂准备一个一般结果. 固定轨道$\ell$, 乘自旋$1/2$共$2(2\ell+1)$维. 最大$m=\ell+1/2$只有$\ket{\ell,\ell}\ket+$, 被$J_+$湮灭, 所以它生成$j_+=\ell+1/2$的$2\ell+2$维阶梯. 其正交补在$J_i$作用下保持, 还剩$2\ell$维. 当$\ell\geq1$, $m=\ell-1/2$有两候选$\ket{\ell,\ell-1}\ket+$和$\ket{\ell,\ell}\ket-$, 已有阶梯占一个, 与它正交的一个最高态必被$J_+$湮灭. 因而补空间生成$j_-=\ell-1/2$的$2\ell$维阶梯, 维数已全用完. $\ell=0$只剩$j=1/2$, 没有负$j$.

具体下支最高态写$a\ket{\ell,\ell-1}\ket++b\ket{\ell,\ell}\ket-$, 条件$J_+=0$给$\sqrt{2\ell}a+b=0$, 归一取$a=1/\sqrt{2\ell+1},b=-\sqrt{2\ell/(2\ell+1)}$. 这同时给末态自旋测量概率, 后面2012期末衰变题直接使用. “最高权态”只指一组中最大$m$的态, 并非另外一种物理粒子.
